% arara: lualatex: { shell: no } until !found('log', '\\(?(R|r)e\\)?run (to get|LaTeX)')
\documentclass[
paper=a4,
DIV=11,
fontsize=12pt,
twoside=true,
parskip=half-,
department=FakS,
logostyle=flat
]{OTHRartcl}


\setlength{\overfullrule}{2pt}

\usepackage[german]{babel}
\usepackage{microtype}
\usepackage{tcolorbox}
\usepackage{booktabs}
\usepackage{tikz}
\usetikzlibrary{shadows, decorations.pathmorphing}
\usepackage[tikz]{bclogo}
\usepackage{colortbl}
\usepackage{tabularx}
\usepackage{latexsym}
\usepackage[hidelinks]{hyperref}
\usepackage{tablefootnote}
\usepackage{csquotes}

\defineUserLogo{OTHR_OTHRLaTeX_Logo}
\shorthandon{"}
\newcommand{\docoption}[1]{\marginpar{\tiny\ttfamily #1}}
\newcommand{\datei}[1]{{\ttfamily #1}}
\usepackage{lipsum}

\newcommand{\Wichtig}[1]{\begin{bclogo}[margeG=0,margeD=0,couleur=OTHR25,
 logo=\bctakecare
]{Wichtig:}#1\end{bclogo}}

%\useDepartmentLogo

\title{Das \LaTeX"=Paket der OTH Regensburg}
\author{OTHR"=\LaTeX"=Team}
\documenttype{Benutzerhandbuch}
\date{\today}

\usepackage{biblatex}
\addbibresource{OTHRtex.bib}


\begin{document}

\maketitle

\tableofcontents
\clearpage

\section{Übergreifende Dokumentoptionen}

\begin{description}
    \item \texttt{department}\docoption{department} Diese Option nimmt diverse Einstellungen entsprechend der jeweils spezifizierten Einheit im Dokument vor. Dies betrifft z.B. Farben.
        Dadurch wird eine einfache Anpassung von Dokumenten möglich.
        Details, inklusive der unterstützten Optionen, finden Sie im Abschnitt~\ref{Sec:Farben} und Beispiele im Anhang~\ref{Sec:Samples}.
        Wird diese Option weggelassen, so wird automatisch der Defaultwert \texttt{OTHR} verwendet.
    \item \texttt{blackFont}\docoption{blackFont} Schaltet die über die Option \texttt{department} umschaltbaren Farben temporär auf Schwarz--Weiß--Betrieb, um eine verbesserte Ausgabe beim Druck zu erhalten.
        Es werden nur die \LaTeX"=Farben umgestellt, die Farben in Grafiken und Logos bleiben ggf. unverändert erhalten.
    \item \texttt{Lucida}\docoption{Lucida} Verwendung der kommerziellen Schrift Lucida als Dokumentschrift.
        Dies benötigt einen Zukauf einer Lucida-Lizenz und Installation des Zeichensatzes..
    \item \texttt{LucidaGrande}\docoption{LucidaGrande}  Verwendung der kommerziellen Schrift LucidaGrande als Dokumentschrift wie sie bei einigen Apple-Systemen vorinstalliert ist.
    \item \texttt{KeepRoman}\docoption{KeepRoman} Für lange geschlossene Texte ist serifenlose Schrift, die per Default gewählt wird, suboptimal in der Lesbarkeit und wirkt insbesondere ohne weitere flankierende typografische Maßnahmen leicht ermüdend.
        Daher kann mit dieser Option die serifenbehaftete Schrift als Standardschrift für das Dokument beibehalten werden.
\end{description}

\section{Vorhandene Dokumententypen}\label{Sec:Klassen}

\begin{center}
\begin{tabularx}{.9\linewidth}{lX}\toprule
\textbf{Dokumententyp} & \textbf{Beschreibung} \\\midrule
\texttt{OTHRartcl} & Dokumentklasse basierend auf der KOMA"=Skript Klasse \texttt{scrartcl} mit Einbindung aller im folgenden beschriebenen Teilpakete zur Farb"=, Font"= und Logowahl, der offiziellen Titelseite, und Zusätzen für die formgerechte Erstellung von Abschlussarbeiten;\\
\texttt{OTHRreprt} & dito, basierend auf der KOMA"=Skript Klasse \texttt{scrreprt};\\
\texttt{OTHRbook} & dito, basierend auf der KOMA"=Skript Klasse \texttt{scrbook};\\
Briefvorlage & Eine mit \texttt{scrlttr2} zu verwendende Stildatei, welche das Layout der OTHR"=Briefvorlage umsetzt;\\
Präsentation & Ein Beamer"=Template zur Erstellung von Präsentationen entsprechend der CI"=Vorgaben;\\
\bottomrule
\end{tabularx}
\end{center}

In den folgenden Unterabschnitten werden lediglich die spezifischen Besonderheiten der jeweiligen Dokumentarten beschrieben.
Alle Vorlagen unterstützen die in den danach folgenden Kapiteln beschriebenen Funktionen zur Schriftart"=, Farb"= und Logoauswahl.

\subsection{\texttt{OTHRartcl}, \texttt{OTHRreprt}, \texttt{OTHRbook}}

Diese Klassen unterstützen eine Titelseite oder auch den titelseitenlosen Modus (Dokumentenoption \verb|notitlepage|).
Die Daten für die Titelseite werden wie bei den Dokumentklassen in KOMA"=Skript üblich mit \verb|\author|, \verb|\title| und ggf. \verb|\date| gesetzt (siehe~\cite{KOMA}) .

Zusätzlich erfordern diese Klassen die Definition eines Dokumenttyps (beliebiger Text wie \enquote{Bericht} oder \enquote{Skript} zur korrekten Darstellung der Titelseite.

Der Titel wird wie gewohnt mit \verb|\maketitle| erzeugt.

\begin{center}
\begin{tabularx}{.9\linewidth}{lX}\toprule
\textbf{Kommando} & \textbf{Erklärung} \\\midrule
\verb+\documenttype{}+ & Definiert den Dokumenttyp, z.B. Masterarbeit, Bachelorarbeit, Projektarbeit, Studienarbeit, Handbuch oder Versuchsanleitung.
Dieser wird für die Titelseite bzw. den Dokumentkopf verwendet.\\
\verb+\studentid{}+ & Zum Setzen der Matrikelnummer auf die Titelseite von Abschlussarbeiten\\
\verb+\department{}+ & Zum Setzen des Namens der Fakultät auf die Titelseite von Abschlussarbeiten\\
\verb+\studyprogramme{}+ & Name des Studiengangs, primär für die Titelseite von Abschlussarbeiten\\
\verb+\startingdate{}+ & Beginn der Arbeit, ebenfalls für Abschlussarbeiten\\
\verb+\closingdate{}+ & Abgabetermin, ebenfalls für Abschlussarbeiten\\
\verb+\firstadvisor{}+ & Erstprüfer, ebenfalls für Abschlussarbeiten\\
\verb+\secondadvisor{}+ & Zweitprüfer, ebenfalls für Abschlussarbeiten\\
\verb+\externaladvisor{}+ & Externe Betreuung , ebenfalls für Abschlussarbeiten\\
\verb+\externallogo[]{}+ & Optionales Logo für die Titelseite; Diese Option ist mit Vorsicht zu genießen und wird ausdrücklich nicht empfohlen. Das optionale Argument wird an \verb|\includegraphcs| weitergereicht.\\
\bottomrule
\end{tabularx}
\end{center}
Außerdem ist es noch möglich, mit \verb|\makedeclaration| eine Erklärung zur selbständigen Bearbeitung der Aufgabe auszugeben.
Die Korrektheit des Textes bezüglich der aktuellen Anforderungen der relevanten  Prüfungsordnungen an diesen Text ist zu prüfen.

\subsection{Briefvorlage}

Die Biefklasse stützt sich auf die KOMA"=Skript Klasse \verb|scrlettr2|.
Für weitere Details beachten Sie bitte die Beispieldatei sowie die Anleitung zu \verb|scrlettr2| in~\cite{KOMA}.


\subsection{Präsentation}

Das Präsentationstemplate wurde in 2025 komplett neu erstellt. Die zugehörige Dokumentation findet sich in \cite{othr-beamer}.

\subsection{Andere Dokumenttypen}

Nicht immer sind die hier vorgestellten Dokumentklassen bzw. Beispieldokumente für die eigene Anwendung passend.
Für diese Fälle kann die entsprechende Infrastruktur zur Erzeugung eines Dokuments nach der CI"=Richtlinie auch über das Paket \datei{OTHR.sty} in ein beliebiges Dokument eingebunden werden.

\begin{verbatim}
\usepackage{OTHR}
\end{verbatim}

Sollten in Spezialfällen nur Teile davon benötigt werden, so können die Pakete \texttt{othr-colors} und \texttt{other-fonts} auch einzeln eingebunden werden.

Damit stehen die ausgewählten vordefinierten Elemente für eigene Dokumente zur Verfügung.

\section{Farben}\label{Sec:Farben}

\subsection{Adaptive Farbe}

Die Farbdefinition \texttt{Accent} und ihre Verwandtschaft passt sich der Dokumentenoption \verb+department+\docoption{department} an und stellt so eine konsistente Farbdarstellung über das Dokument sicher - auf beim Wechsel der \texttt{department}"=Einstellung.

Hier sind die Farbnamen zur Verwendung der adaptiver Farben am Beispiel der gewählten Stiloption \verb+department=FakSG+ dargestellt.

\verb|\textcolor{Accent}{...}|\\
\textcolor{Accent}{\bfseries Die Akzentfarbe \texttt{Accent} wird automatisch an die Fakultät angepasst.}

\verb|\textcolor{Accent90}{...}|\\
\textcolor{Accent90}{\bfseries Die Akzentfarbe \texttt{Accent90} wird automatisch an die Fakultät angepasst.}

\verb|\textcolor{Accent50}{...}|\\
\textcolor{Accent50}{\bfseries Die Akzentfarbe \texttt{Accent50} wird automatisch an die Fakultät angepasst.}

\verb|\textcolor{Accent25}{...}|\\
\textcolor{Accent25}{\bfseries Die Akzentfarbe \texttt{Accent25} wird automatisch an die Fakultät angepasst.}

\verb|\textcolor{AccentFontColor}{...}|\\
\textcolor{AccentFontColor}{Die Akzentfarbe \texttt{AccentFontColor} wird automatisch an die Fakultät angepasst, wobei bei im unterschied zur Farbe \texttt{Accent} bei Farben mit mangelndem Kontrast vor weißem Hintergrund  auf \texttt{OTHgrau} ausgewichen wird.}

\begin{center}
\colorbox{Accent}{\begin{minipage}{.9\linewidth}\color{AccentContrastColor} Die Farbdefinition \texttt{AccentContrastColor} kann verwendet werden, um sicher lesbaren Text auf einen Hintergrund in Akzentfarbe zu setzen.\end{minipage}}
\end{center}

Die Dokumentenoption \verb+blackFont+ \docoption{blackFont} setzt die \texttt{AccentFontColor} ungeachtet der verwendeten Stiloption \verb+department+ auf schwarz und die \texttt{AccentContrastColor} auf weiß.
Damit lässt sich bei Nutzung der farbigen Textgestaltung z.B. in Präsentationen leicht auf ein evtl. druckfreundlicheres Farbschema wechseln.

\Wichtig{\sloppy Es wird dringend empfohlen, die vorhandenen adaptiven Farben \texttt{Accent}, \texttt{Accent90}, \texttt{Accent50}, \texttt{Accent25}, \texttt{AccentFontColor} und bei Bedarf auch \texttt{AccentContrastColor} zu verwenden.
Dadurch passt sich die Gestaltung des Dokuments automatisch korrekt an die Dokumentenoption \texttt{department} an.}


Derzeit unterstützte Werte für die Option \verb+department+\docoption{department} und ihre Bedeutung:

\begin{table}[!h!]
\centering\scriptsize
\begin{tabular}{lcccc}\toprule
\textbf{Wert} & Einrichtung & Farbschema & \verb+AccentFontColor+ & \verb+AccentContrastColor+\\\midrule
OTHR & alle & Schwarz/grau & \verb+OTHRgrau+ & \verb+black+ \\
FakA & Fakultät A & Fakultätsfarbe & \verb+OTHRgrau+ & \verb+black+ \\
%\textcolor{black!20}{FakAM} & \textcolor{black!20}{Fakultät AM} & \textcolor{black!20}{Fakultätsfarbe} & \textcolor{black!20}{Fakultätsfarbe} & \color{black!20}\verb+white+ \\
FakANK & Fakultät ANK & Fakultätsfarbe & Fakultätsfarbe & \verb+white+ \\
FakB & Fakultät B & Fakultätsfarbe & Fakultätsfarbe & \verb+white+ \\
FakBM & Fakultät BM & Fakultätsfarbe & Fakultätsfarbe & \verb+white+ \\
FakEI & Fakultät EI & Fakultätsfarbe & Fakultätsfarbe & \verb+white+ \\
FakIM & Fakultät IM & Fakultätsfarbe & Fakultätsfarbe & \verb+white+ \\
FakM & Fakultät M & Fakultätsfarbe & Fakultätsfarbe & \verb+white+ \\
FakSG & Fakultät SG & Fakultätsfarbe & Fakultätsfarbe & \verb+white+ \\
ZWW & ZWW & ZWW-Farbe & \verb+OTHgrau+ & \verb+black+ \\
Labornamen & Labore & diverse & diverse & diverse \\
\bottomrule
\end{tabular}
\end{table}

Hierbei ist die Farbe \verb+AccentFontColor+ jeweils für Text auf weißem Grund gedacht und wird im Allgemeinen auf die Farbe des mit der Stiloption  \verb+department+ gewählten Bereichs gesetzt.
Da bei manchen der vordefinierten Farben aber der Kontrast auf weißem Hintergrund zu schwach ausfällt, wird in diesen Fällen Hochschulgrau für \verb+AccentFontColor+ gewählt.

Zudem zeigt sich, dass bei Verwendung der Bereichsfarben als Hintergrund je nach Farbe der Einrichtung eine andere Schriftfarbe nötig ist, um ausreichenden Kontrast zu erreichen.
Durch Verwendung der Farbe \verb+AccentContrastColor+ kann hier eine automatische Anpassung im Dokument vorgenommen werden.

\subsection{Farbübersicht}

Auch die Farben der einzelnen Einheiten stehen als stabile Definitionen unabhängig von der Stiloption \verb+department+\docoption{department} zur Verfügung.

\newcommand{\farbTabellenZeile}[1]{
{\color{#1}\texttt{#1}} &
{\color{#190}\texttt{#190}} &
{\color{#150}\texttt{#150}} &
{\color{#125}\texttt{#125}} &
{\color{#1FontColor}\texttt{#1FontColor}} &
{\cellcolor{#1}\color{#1ContrastColor}\texttt{#1ContrastColor}} \\
}
\begin{center}\scriptsize
\begin{flushleft}\small\textbf{Fakultäten und Einrichtungen}\end{flushleft}\nopagebreak
\begin{tabular}{cccccc}
\farbTabellenZeile{Accent}
\farbTabellenZeile{OTHR}
\farbTabellenZeile{FakA}
\farbTabellenZeile{FakANK}
\farbTabellenZeile{FakAM}
\farbTabellenZeile{FakB}
\farbTabellenZeile{FakBM}
\farbTabellenZeile{FakBW}
\farbTabellenZeile{FakEI}
\farbTabellenZeile{FakIM}
\farbTabellenZeile{FakM}
\farbTabellenZeile{FakSG}
\farbTabellenZeile{FakS}
\farbTabellenZeile{ZWW}
\\
\farbTabellenZeile{RCAI}
\farbTabellenZeile{RCER}
\farbTabellenZeile{RCHST}
\farbTabellenZeile{RSDS}
\\
\end{tabular}

\vspace{\fill}\pagebreak

\begin{flushleft}\small\textbf{Laborfarben}\end{flushleft}\nopagebreak
\begin{tabular}{cccccc}
%A
\farbTabellenZeile{FMIS}
%ANK
\farbTabellenZeile{NACH}
\farbTabellenZeile{SappZ}
%B
\farbTabellenZeile{Geo}
\farbTabellenZeile{KNB}
%BW
\farbTabellenZeile{12Science}
\farbTabellenZeile{FEURO}
%EI
\farbTabellenZeile{Bisp}
\farbTabellenZeile{DK0PT}
\farbTabellenZeile{FENES}
\farbTabellenZeile{LAS3}
\farbTabellenZeile{MRU}
\farbTabellenZeile{SES}
%IM
\farbTabellenZeile{CCSE}
\farbTabellenZeile{eHealth}
\farbTabellenZeile{IMDS}
\farbTabellenZeile{LFD}
\farbTabellenZeile{MD}
\farbTabellenZeile{ReMIC}
\farbTabellenZeile{SEC}
%M
\farbTabellenZeile{AIMS}
\farbTabellenZeile{BFM}
\farbTabellenZeile{BMA}
\farbTabellenZeile{CEEC}
\farbTabellenZeile{CFD}
\farbTabellenZeile{FEM}
\farbTabellenZeile{IPF}
\farbTabellenZeile{KIB}
\farbTabellenZeile{KWK}
\farbTabellenZeile{LAT}
\farbTabellenZeile{LBM}
\farbTabellenZeile{LeanLab}
\farbTabellenZeile{LFT}
\farbTabellenZeile{LFW}
\farbTabellenZeile{LMP}
\farbTabellenZeile{LMS}
\farbTabellenZeile{LRT}
\farbTabellenZeile{LWS}
\farbTabellenZeile{LWM}
\farbTabellenZeile{MKS}
\farbTabellenZeile{MST}
\farbTabellenZeile{RRRU}
\farbTabellenZeile{RST}
%
\farbTabellenZeile{IST}
\farbTabellenZeile{LP}
\farbTabellenZeile{PF}
\farbTabellenZeile{PT}


\end{tabular}
\end{center}

\section{Schriftarten: \datei{OTHR\_fonts.sty}}

Die Hausschrift des Corporate Designs der OTH Regensburg ist die Lucida Sans.
Um diesem Schriftbild möglichst nah zu kommen, unterstützt das Zeichensatzpaket verschiedene Optionen.

Default ist die Verwendung der DejaVU Sans, welche dem Original sehr ähnlich ist.
Diese Schrift wird bei vielen Distributionen im Standardumfang mitgeliefert, bietet aber in der Type~1 Variante keinen Zeichensatz für den Formelsatz, der daher mit pdflatex traditionell gesetzt wird.

Bei Verwendung von LuaTeX wird versucht, als Standard eine der Originalschriftarten einzusetzen, sofern diese Auf dem Rechner installiert sind.
Alternativ wird auch hier die OpenType"=Variante von DejaVU Sans inklusive des Zeichensatzes für den Formelsatz verwendet.

Falls eine Lizenz für die Lucida"=Zeichensätze vorhanden ist, kann auch bei pdflatex über die Option \texttt{Lucida}\docoption{Lucida\\lucida} auf die Lucida"=Schriftfamilie umgeschaltet werden, die auch den Formelsatz unterstützt und insgesamt ein sehr harmonisches Schriftbild liefert\footnote{Lizenzen für die Lucida Fonts im OpenType"= und T1"= Format sind bei der \TeX Users Group erhältlich: \href{https://www.tug.org/store/lucida/index.html}{https://www.tug.org/store/lucida/index.html}.}.

Bei Verwendung von LuaTeX und einer vorhandenen Installation der Lucida Schriftart Grande auf dem System (typisch für Apple"=Systeme), kann über die Option \texttt{LucidaGrande}\docoption{LucidaGrande} auch diese Schriftart verwendet werden.

Da sich in manchen Fällen die serifenlose Schriftart nur bedingt eignet (z.B. sehr lange geschlossene Texte mit relativ langen Zeilen bzw. hoher Zeichenanzahl) kann über die Option \texttt{KeepRoman}\docoption{KeepRoman} die Variante der jeweils gewählten Schrift mit Serifen als Standardschrift beibehalten werden.


\end{document}
